\section*{第 \ref{ch:g7:areas} 章——面积与体积}

\begin{solution}{exo:g7:areas:1}
$8 \times 5 = 40$ cm$^2$；\qquad $6.5 \times 4 = 26$ m$^2$。
\end{solution}

\begin{solution}{exo:g7:areas:2}
$A = 10 \times 4 = 40$ cm$^2$。答案不是
$10 \times 6 = 60$，因为 $6$ cm 的那条边是倾斜的：公式
用的是\emph{高}（\omterm{def:g6:lines:perp}{垂直}距离），它是 $4$ cm。
\end{solution}

\begin{solution}{exo:g7:areas:3}
$\frac{12 \times 7}{2} = 42$ cm$^2$; \quad
$\frac{9 \times 6}{2} = 27$; \quad
$\frac{5.5 \times 2}{2} = 5.5$ m$^2$。
\end{solution}

\begin{solution}{exo:g7:areas:4}
圆盘：$\pi \times 5^2 = 25\pi \approx 79$ cm$^2$。\omterm{def:g3:shapes:shapes}{半径}为
$10$ 的半圆：$\frac{\pi \times 10^2}{2} = 50\pi \approx 157$ cm$^2$。
\end{solution}

\begin{solution}{exo:g7:areas:5}
\omterm{def:g6:measure:perimeter}{周长}：$2\pi \times 3 = 6\pi \approx 18.8$ cm（长度，单位为 cm）。
\omterm{def:g6:measure:area}{面积}：$\pi \times 3^2 = 9\pi \approx 28.3$ cm$^2$（\omterm{def:g6:measure:area}{面积}，单位为
cm$^2$）。
\end{solution}

\begin{solution}{exo:g7:areas:6}
\omterm{def:g7:areas:prism}{棱柱}：$24 \times 10 = 240$ cm$^3$。\omterm{def:g7:areas:prism}{圆柱}：
$\pi \times 3^2 \times 8 = 72\pi \approx 226$ cm$^3$。
\end{solution}

\begin{solution}{exo:g7:areas:7}
$\frac{8 \times h}{2} = 28$，于是 $4h = 28$，从而 $h = 7$ cm。
\end{solution}

\begin{solution}{exo:g7:areas:8}
三个乘积相等，因为每一个都算的是\emph{同一个}
量——三角形的\omterm{def:g6:measure:area}{面积}。这给出了一种实用的
作图检验，也是一种由另一组
底--高对求高的方法（用于\cref{exo:g7:areas:12}）。
\end{solution}

\begin{solution}{exo:g7:areas:9}
两个三角形的\omterm{def:g6:measure:area}{面积}分别为 $\frac{30 \times 20}{2} = 300$ m$^2$
和 $\frac{18 \times 20}{2} = 180$ m$^2$，合计 $480$ m$^2$。由此可猜出
一般公式：对于\omterm{def:g6:lines:perp}{平行}边 $a$ 和 $b$，它们之间的距离为 $h$，
\[
A = \frac{(a + b) \times h}{2}
\]
——梯形\omterm{def:g6:measure:area}{面积}公式（此处为 $\frac{(30+18)\times 20}{2} = 480$）。
\end{solution}

\begin{solution}{exo:g7:areas:10}
水的\omterm{def:g6:measure:volume}{体积}：$\pi \times 6^2 \times 15 = 540\pi \approx 1\,696$
cm$^3$。\omterm{def:g5:solids:def}{长方体}底\omterm{def:g6:measure:area}{面积}：$18 \times 12 = 216$ cm$^2$。达到的高度：
$\frac{540\pi}{216} = 2.5\pi \approx 7.9$ cm。
\end{solution}

\begin{solution}{exo:g7:areas:11}
\omterm{def:g3:shapes:shapes}{半径}分别为 $15$ 和 $20$ cm。\omterm{def:g6:measure:area}{面积}：$225\pi \approx 707$ cm$^2$ 和
$400\pi \approx 1257$ cm$^2$。每 $100$ cm$^2$ 的价格：
$\frac{9}{7.07} \approx 1.27$ 欧元，
$\frac{14}{12.57} \approx 1.11$ 欧元。大号比萨更划算
——\omterm{def:g6:measure:area}{面积}随\omterm{def:g4:geometry:circle}{直径}的\emph{平方}增长。
\end{solution}

\begin{solution}{exo:g7:areas:12}
以两直角边为底与高：$A = \frac{6 \times 8}{2} = 24$ cm$^2$。以
斜边为底：$24 = \frac{10 \times h}{2} = 5h$，于是
$h = 4.8$ cm。
\end{solution}

\begin{solution}{pb:g7:areas:1}
\textbf{1.} 正方形：$8 \times 8 = 64$。长方形：
$5 \times 13 = 65$。同样的四块拼片在一幅图中似乎覆盖了 $64$
个平方单位，而在另一幅中覆盖了 $65$ 个平方单位：多出了一个单位的
\omterm{def:g6:measure:area}{面积}凭空产生——据称如此。

\textbf{2.} 每个三角形：$\frac{8 \times 3}{2} = 12$。每个
梯形：$\frac{(5 + 3) \times 5}{2} = 20$。

\textbf{3.} $12 + 12 + 20 + 20 = 64$：四块拼片合计 $64$，
因此它们恰好能铺满正方形——却\emph{不能}铺满
\omterm{def:g6:measure:area}{面积}为 $65$ 的长方形。这幅长方形图中必定含有一个
空隙。

\textbf{4.} 若坡度相同，则 $3$ 比 $8$
应与 $2$ 比 $5$ 成比例；交叉相乘得：
$3 \times 5 = 15$ 与 $2 \times 8 = 16$。两者不相等：
斜边与倾斜的边\emph{并不}对齐。长方形的
“对角线”是一个谎言——两条斜率略有不同的
线在一个浅角处相交。

\textbf{5.} 沿着这条假的对角线，四块拼片留下了一条
又长又极细的缝隙——形状像一个非常
扁平的\omterm{def:g7:central:parallelogram}{平行四边形}——其\omterm{def:g6:measure:area}{面积}恰好是缺少的
$65 - 64 = 1$ 个平方单位。沿一条约
$13$ 个单位长的对角线展开，其平均宽度约为
$1 \div 13 \approx 0.08$ 个单位：在真实的图上，它比
遮住它的铅笔线还要细。

\textbf{6.} 每个三角形的底为 $6$ cm，高为 $4$ cm——
两\omterm{def:g6:lines:perp}{平行}线之间的距离，无论\omterm{def:g5:solids:def}{顶点}位于何处——
因此每个\omterm{def:g6:measure:area}{面积}为 $\frac{6 \times 4}{2} = 12$ cm$^2$
（见\cref{thm:g7:areas:triangle}）。沿着\omterm{def:g6:lines:perp}{平行}线移动\omterm{def:g5:solids:def}{顶点}
只改变形状，底或高从不改变：
\omterm{def:g6:measure:area}{面积}永远相等。

\textbf{7.} 对角线把梯形分成两个三角形，
一个以 $B$ 为底，另一个以 $b$ 为底，两者的高都是 $h$（
\omterm{def:g6:lines:perp}{平行}两边之间的距离）：
\[
A = \frac{B \times h}{2} + \frac{b \times h}{2}
= \frac{(B + b) \times h}{2} .
\]

\textbf{8.} $A = \frac{(9 + 5) \times 4}{2} = \frac{56}{2} =
28$ cm$^2$.

\textbf{9.} \omterm{def:g6:lines:perp}{平行}两边的平均数是
$\frac{B + b}{2}$，于是（平均数）$\times h =
\frac{(B+b) \times h}{2}$：同一公式，只是因式分解
的方式不同。检验：平均数 $= \frac{9+5}{2} = 7$，而
$7 \times 4 = 28$ cm$^2$。

\textbf{10.} 斜放的这一摞卡片，所含卡片与原来完全相同
——层数相同，每层\omterm{def:g6:measure:area}{面积}也相同——因此
\omterm{def:g6:measure:volume}{体积}也相同。在 $V = B \times h$ 中，高必须是与底面
\emph{\omterm{def:g6:lines:perp}{垂直}}的高（即这一摞卡片的竖直高度）
，而不是沿斜边量取：正如\omterm{def:g7:central:parallelogram}{平行四边形}的
\omterm{def:g6:measure:area}{面积}只需要\omterm{def:g6:lines:perp}{垂直}高，而非斜边长
（见\cref{exo:g7:areas:2}），只是升高了一维。

\textbf{11.} 圆环\omterm{def:g6:measure:area}{面积} $= \pi \times 5^2 - \pi \times 3^2 = 25\pi -
9\pi = 16\pi \approx 50$ m$^2$（见\cref{thm:g7:areas:disk}）。

\textbf{12.} \omterm{def:g6:measure:area}{面积}为 $16\pi$ 的圆，\omterm{def:g3:shapes:shapes}{半径}为 $4$，因为
$\pi \times 4^2 = 16\pi$。数 $3$、$4$、$5$ 是数学中最著名的
一组勾股数，将在
\cref{ch:g8:pythagoras}中担任主角。

\textbf{13.} 一个\omterm{def:g4:geometry:circle}{直径} $40$ cm 的披萨：\omterm{def:g3:shapes:shapes}{半径} $20$，
\omterm{def:g6:measure:area}{面积} $\pi \times 400 \approx 1\,257$ cm$^2$。两个 $28$ cm
的披萨：\omterm{def:g3:shapes:shapes}{半径} $14$，合计
$2 \times \pi \times 196 = 392\pi \approx 1\,232$ cm$^2$。
单个大披萨的量比两个中号披萨\emph{多}。

\textbf{14.} 将\omterm{def:g3:shapes:shapes}{半径}扩大为 $1.4$ 倍，则\omterm{def:g6:measure:area}{面积}扩大为
$1.4^2 = 1.96$ 倍——几乎是两倍。精确的\omterm{def:g5:division:multiple}{倍数}是
平方为 $2$ 的数，约为 $1.414$：正是
\cref{ex:g8:pythagoras:diagonal}中的对角线数。（经验法则：
\omterm{def:g4:geometry:circle}{直径}增加\omterm{def:g3:division:half}{一半}，披萨几乎是两倍。）

\textbf{15.} $13 \times 13 = 169$，但 $8 \times 21 = 168$：
这一次有一个平方单位\emph{消失}了——各块
拼合处沿一条细长条重叠，而不是留下空隙。
取赫马昌德拉–斐波那契数列中的连续项 $3, 5, 8, 13, 21$，
乘积在比平方数多 1
与少 1 之间交替（$5 \times 13 = 65 =
8^2 + 1$，随后 $8 \times 21 = 168 = 13^2 - 1$）：魔术师
交替地“多出”和“少了”一个单位。\omterm{def:g6:measure:area}{面积}从不说谎：
各块的总\omterm{def:g6:measure:area}{面积}为 $64$，无论怎样滑动，都覆盖 $64$——每一个
缺少或多出的一个单位，都藏在一条细长条里，
等待一次斜率检验将其揭示出来。
\end{solution}
